\documentclass[10pt,a4paper]{amsart}
\usepackage[margin=19mm]{geometry}
\usepackage{amsmath,amssymb,mathtools}
\usepackage[scr=rsfs,cal=euler]{mathalfa}
\usepackage{xfrac}
\usepackage[unicode,hidelinks,bookmarks=false]{hyperref}
\newcommand{\ie}{i.e.\ }
\newcommand{\cf}{cf.\ }
\newcommand{\resp}{respectively\ }
\newcommand{\Subsection}{\subsection}
\providecommand{\nobreakdash}{\nobreak\hbox{-}}
\setlength{\emergencystretch}{2em}
\multlinegap0pt
\usepackage{enumerate}
\usepackage{amssymb}
\usepackage{float}
\usepackage{bm}
\usepackage{algorithm}\usepackage{algpseudocode}
\newtheorem{definition}{Definition}
\newtheorem{assumption}{Assumption}
\newtheorem{lemma}{Lemma}
\def\bv{\mathbf{v}}
\title{\bf {Bias-Corrected Multiplier Bootstrap Inference for Spectral Edges of Large Covariance Matrices}}
\author{Xiucai Ding, Yichen Hu, Jiahui Xie}
\date{Source: arXiv:2607.08089v2 (CC BY 4.0)}
\begin{document}
\maketitle

\noindent\emph{Source and licence.} \bf {Bias-Corrected Multiplier Bootstrap Inference for Spectral Edges of Large Covariance Matrices}. Version arXiv:2607.08089v2 (CC BY 4.0), distributed by arXiv under the Creative Commons Attribution 4.0 International licence, http://creativecommons.org/licenses/by/4.0/, as recorded in the arXiv licence metadata for this exact version and shown on the abstract page https://arxiv.org/abs/2607.08089v2. Full licence notice: \texttt{licenses/arxiv-2607.08089-CC-BY-4.0.txt}.\par\smallskip

\noindent\textit{Editorial scope.} Counted unit: the article's lemma lem\_stabilityargument (source0.tex lines 2758--2771). The article's complete original proof of it is delivered with it (source0.tex lines 2775--2905, one \texttt{proof} environment of 131 lines, which the article places away from the statement). No mathematics is written, repaired or re-derived: the statement and the proof are the article's own lines, in the article's own notation, and no retained line is substituted or re-flowed.  Retained uncounted as the source-local context the proof needs: lines 543--562; lines 833--836; lines 1859--1865; lines 1930--1939; lines 1951--1978; lines 2622--2651; lines 2671--2685; lines 2749--2756.  Nothing was omitted from any retained span, so the package carries no omission record.  No retained line contains a citation, so the wrapper carries no bibliography and no bibliography entry.  No retained line contains a figure environment, an image-inclusion command or drawing code, so no image is delivered and the package declares no asset dependency.  Editorial formatting only, in addition: an AMS \texttt{amsart} preamble that restates the article's own declarations and macros used by the retained lines, namely the environments \texttt{definition}, \texttt{assumption}, \texttt{lemma} on separate counters, together with the standard packages those lines need.  The retained environments print in this document as Definition 1, Assumption 1, Lemma 1 in order of appearance, whereas the article prints the same items with the numbering of its own full document.  The complete native source of the article as distributed in the arXiv e-print for this exact version is bundled unchanged as \texttt{sources/204756/source0.tex}.
\begin{definition}\label{defn_feasiblemultiplier}
	\normalfont
	Throughout the paper, we call a multiplier $\xi^2$ \emph{feasible} if it satisfies the following conditions:
	\begin{enumerate}
		%	    \setlength{\itemsep}{0.3em}
		%	\setlength{\parsep}{0pt}
		%	\setlength{\topsep}{0.3em}
		%	\setlength{\labelsep}{0.6em}
		\item Its mean is normalized as
		\begin{equation}\label{eq_meancondition}
			\mathbb{E}\xi^2 = 1.
		\end{equation}
		
		\item  Its variance satisfies
		\begin{equation}\label{eq_variancecondition}
			\operatorname{Var}(\xi^2) \asymp n^{-1/3+\delta},
		\end{equation}
		for some constant $\delta$ such that $2\delta_* < \delta < 1/3$, where $\delta_*>0$ is a sufficiently small constant to be specified later in Definition~\ref{def_OmegaX} of the supplement.
	\end{enumerate}
\end{definition}

\begin{equation}\label{eq_basepopulationcovariancematrix}
	\vspace*{-10pt}
	\Sigma_0=\sum_{i=1}^p \sigma_i \bv_i \bv_i^\top. 
\end{equation}

\begin{assumption}\label{assum_technique}
	When multipliers satisfy Definition \ref{defn_feasiblemultiplier}, for $\Sigma_0$ defined in \eqref{eq_basepopulationcovariancematrix}, we assume for some constant $\tau_0>0$ such that 
	\begin{align*}
		\min_{1\le i\le p}|1+\sigma_i m_{2n,c}(\mathsf{E}_{\mathrm{MB}})|\ge\tau_0,\quad \inf_{t\in\operatorname{supp}(F_{\xi^2})}|1+tm_{1n,c}(\mathsf{E}_{\mathrm{MB}})|\ge\tau_0,
	\end{align*}
	where recall that $\mathsf{E}_{\mathrm{MB}}$ is the rightmost edge of the limiting spectral density of $Q_{\mathrm{MB}}^0$.
\end{assumption}

\begin{definition}\label{def_OmegaM}
	Let $C_{\Xi,1}, C_{\Xi,2}, C_{\Xi,3}$ and $C_{\Xi,4}$  be some positive constants and $0<c_{\Xi,1},c_{\Xi,2}, c_{\Xi,3}<1$ are some sufficiently small constants. Denote $\Omega_\Xi\equiv\Omega_{n,\Xi}$ be the event on $\{\xi^2_i\}$ so that the following conditions hold:
	\begin{align*}
		&\frac{1}{n}\sum_{i=1}^n\xi^2_i\leq C_{\Xi,1},\\
		&\left|\frac{1}{n}\sum_{i=1}^n\frac{\xi^2_i}{1+\xi^2_im_{1n,c}(\mathsf{E}_{\mathrm{MB}})}-\int\frac{t}{1+tm_{1n,c}(\mathsf{E}_{\mathrm{MB}})}\mathrm{d}F_{\xi^2}(t)\right|\leq \frac{C_{\Xi,2}n^{c_{\Xi,1}}}{\sqrt{n}},\\
		&\left|\sum_{i=1}^n(\xi^2_i-1)\right|\leq C_{\Xi,3}\sqrt{(n\log^{c_{\Xi,2}}n)\operatorname{Var}\xi^2},\\
		&\frac{1}{n}\sum_{i=1}^n(\xi^2_i-1)^2\leq C_{\Xi,4}\log^{c_{\Xi,3}}n\operatorname{Var}\xi^2.
		%&\max_{1\le i\le n}|\xi^2_i-1|\leq C_{\Xi,5}\log^{-c_{\Xi,4}}n.
	\end{align*}
\end{definition}

\begin{definition}\label{def_OmegaX}
	Denote $\Omega_X$ as the event on $X$ so that the following conditions hold for some small constants $\delta_*$ and $\tau_*$:
	\begin{enumerate}
		\setlength{\itemsep}{0.3em}
		\setlength{\parsep}{0pt}
		\setlength{\topsep}{0.3em}
		\setlength{\labelsep}{0.6em}
		\item [(1)] For any deterministic $p\times p$ matrix $A$, we have all $1\leq i,j\leq n$, the columns of $X$ satisfy
		\begin{align*}
			|\mathbf{x}_i^\top \mathbf{x}_j|=\mathrm{O}\left(n^{c_{X,1}}\sqrt{\frac{\|\mathbf{x}_i\|^2}{n}}\right), \, |\mathbf{x}_i^\top A\mathbf{x}_j|=\mathrm{O}\left(n^{c_{X,1}}\frac{\|A\|_F}{n}\right),\,|\mathbf{x}_i^\top A\mathbf{x}_i-\frac{1}{n}\operatorname{Tr}A|=\mathrm{O}\left(n^{c_{X,1}}\frac{\|A\|_F}{n}\right),
		\end{align*}
		for an arbitrarily small constant $c_{X,1}>0$.
		\item [(2)] $\|XX^\top\|\leq C_{X,1}$ for some constant $C_{X,1}>0$.
		\item [(3)] The eigenvalue rigidity and level repulsion estimates for $Q_{\mathrm{Sam}}^0$ take the form
		\begin{align*}
			&|\mu^0_i-\gamma_i|\leq C_{X,2}(i\wedge n+1-i)^{-1/3}n^{-2/3+c_{X,2}},\quad \text{where}\; n\int_{\gamma_i}^{\infty}\mathrm{d}\varrho=i-\frac{1}{2},\quad i=1,\dots,n;\\
			&|\mu^0_k-\mu^0_{k+1}|\geq C_{X,3}k^{-1/3}n^{-2/3-c_{X,3}}, \quad k=1,\dots, \lfloor \omega p \rfloor;\\
			&|\mu^0_1-\mathsf{E}|\geq C_{X,4} n^{-2/3-c_{X,3}},
		\end{align*}
		for some values $0<\omega<1$, $C_{X,2},C_{X,3},C_{X,4}>0$ and arbitrarily small constants $0<c_{X,2},c_{X,3}\le\delta_*$.
		\item [(4)] For any deterministic diagonal matrix $\Xi^2$, the rigidity of spiked eigenvalues of $Q_{\mathrm{MB}}$ holds as 
		\begin{align*}
			|\lambda_i-\vartheta^{\mathtt{MB}}_i|\le C_{X,5}n^{-1/2+c_{X,4}}\tau_{*}^{1/2},\quad i=1,\dots,r,
		\end{align*}
		where $1+\widetilde{\sigma}_im_{2n}(\vartheta^{\mathtt{MB}}_i)=0$ and $\tau_{*}=\widetilde{\sigma}_i+m_{2n}^{-1}(\widehat{\mathsf{E}}_{\mathrm{MB}})$,
		for some value $C_{X,5}>0$ and arbitrarily small constant $c_{X,4}>0$.
	\end{enumerate}
\end{definition}

\begin{lemma}[Multiplier-side regulation]\label{lem_regulation}
	Let $x_0:=m_{1n,c}(\mathsf{E}_{\mathrm{MB}})$. Suppose Assumption \ref{assum_technique} holds. In particular, assume that there exists a constant $\tau_0>0$ such that
	\begin{align*}
		\inf_{t\in \operatorname{supp}(F_{\xi^2})}
		|1+t x_0|\geq \tau_0 .
	\end{align*}
	Assume further that the multiplier distribution has bounded support, namely
	\begin{align*}
		\operatorname{supp}(F_{\xi^2})\subset [0,\mathsf{C}_\xi]
	\end{align*}
	for some constant $\mathsf{C}_\xi$. Then there exist constants $\delta_0>0$ and $\tau>0$ such that
	\begin{align*}
		\inf_{|x-x_0|\leq \delta_0}
		\inf_{t\in \operatorname{supp}(F_{\xi^2})}
		|1+tx|\geq \tau .
	\end{align*}
	Consequently, for any realization $\{\xi_j^2\}_{j=1}^n$ drawn from
	$F_{\xi^2}$,
	\begin{align*}
		\inf_{|x-x_0|\le \delta_0} \min_{1\leq j\leq n} |1+\xi_j^2x|\geq \tau.
	\end{align*}
	In particular, if
	\begin{align*}
		|m_{1n}(\widehat{\mathsf{E}}_{\mathrm{MB}})-x_0|=\mathrm{o}(1),
	\end{align*}
	then, for all sufficiently large $n$,
	\begin{align*}
		\min_{1\le j\le n}|1+\xi_j^2m_{1n}(\widehat{\mathsf{E}}_{\mathrm{MB}})|\geq \tau .
	\end{align*}
\end{lemma}

\begin{lemma}[Uniform concentration of multiplier transforms]
	\label{lem_uniform_multiplier_concentration}
	Under the assumptions of Lemma \ref{lem_regulation}, for any  small constant
	$c_{\Xi,1}>0$, with probability $1-\mathrm{o}(1)$, the following estimates hold:
	\begin{align*}
		\sup_{|x-x_0|\le\delta_0}\left|\frac{1}{n}\sum_{j=1}^n \frac{\xi_j^2}{1+\xi_j^2x}-\int\frac{t}{1+tx}\,\mathrm{d}F_{\xi^2}(t) \right|\leq n^{-1/2+c_{\Xi,1}},
	\end{align*}
	\begin{align*}
		\sup_{|x-x_0|\le\delta_0}\left|\frac{1}{n}\sum_{j=1}^n\frac{\xi_j^4}{(1+\xi_j^2x)^2}-\int\frac{t^2}{(1+tx)^2}\,\mathrm{d}F_{\xi^2}(t)\right|\leq n^{-1/2+c_{\Xi,1}},
	\end{align*}
	and
	\begin{align*}
		\sup_{|x-x_0|\le\delta_0}\left|\frac{1}{n}\sum_{j=1}^n \frac{\xi_j^6}{(1+\xi_j^2x)^3}-\int\frac{t^3}{(1+tx)^3}\,\mathrm{d}F_{\xi^2}(t)\right|\leq n^{-1/2+c_{\Xi,1}}.
	\end{align*}
\end{lemma}

\begin{definition}[Regular-edge condition]\label{def_regularedge}
	Let $x_0:=m_{1n,c}(\mathsf E_{\mathrm{MB}})$ and $y_0:=\mathsf E_{\mathrm{MB}}$. We say that the edge system is regular if there exists a constant $c_0>0$ such that
	\begin{align}\label{eq_regular_edge_condition}
		|\partial_yF_{n,c}(x_0,y_0)|\ge c_0,
		\qquad
		|\partial_{xx}F_{n,c}(x_0,y_0)|\ge c_0.
	\end{align}
\end{definition}

\begin{lemma}[Stability of the edge system]
	\label{lem_stabilityargument}
	Let $x_0:=m_{1n,c}(\mathsf E_{\mathrm{MB}})$ and $y_0:=\mathsf E_{\mathrm{MB}}$. Suppose that the system is regular in the sense that Definition \ref{def_regularedge} holds.
	Then there exists a unique local pair $(x_1,y_1)$ such that
	\begin{align*}
		F_n(x_1,y_1)=0,
		\qquad
		\partial_xF_n(x_1,y_1)=0
	\end{align*}
	and
	\begin{align*}
		|x_1-x_0|+|y_1-y_0|=\mathrm{O}(n^{-1/2+c_{\Xi,1}}).
	\end{align*}
\end{lemma}

\begin{proof}
	Fix a realization of multipliers in $\Omega_{\Xi}$. The following argument is deterministic.
	Define the two-dimensional maps
	
	\begin{align*}
		G_n(x,y):=
		\begin{pmatrix}
			F_n(x,y)\\
			\partial_xF_n(x,y)
		\end{pmatrix},
		\qquad
		G_{n,c}(x,y):=
		\begin{pmatrix}
			F_{n,c}(x,y)\\
			\partial_xF_{n,c}(x,y)
		\end{pmatrix}.
	\end{align*}
	By the definition of $(x_0,y_0)$ in Lemma \ref{lem_stabilityargument}, we have $G_{n,c}(x_0,y_0)=0.$ Moreover,
	\begin{align*}
		DG_{n,c}(x_0,y_0)
		=
		\begin{pmatrix}
			\partial_xF_{n,c}(x_0,y_0)
			&
			\partial_yF_{n,c}(x_0,y_0)\\
			\partial_{xx}F_{n,c}(x_0,y_0)
			&
			\partial_{xy}F_{n,c}(x_0,y_0)
		\end{pmatrix}.
	\end{align*}
	Since $\partial_xF_{n,c}(x_0,y_0)=0,$
	we have
	\begin{align*}
		\det DG_{n,c}(x_0,y_0)
		=
		-\partial_yF_{n,c}(x_0,y_0)
		\partial_{xx}F_{n,c}(x_0,y_0).
	\end{align*}
	Therefore, by the non-degeneracy assumption,
	\begin{align*}
		|\det DG_{n,c}(x_0,y_0)|\ge c_0^2.
	\end{align*}
	Thus $DG_{n,c}(x_0,y_0)$ is invertible and its inverse is uniformly bounded.
	
	Let $\mathcal N_{\delta_0}:=\{(x,y): |x-x_0|+|y-y_0|\le \delta_0\}$ for some sufficiently small constant $\delta_0>0$. By Definition \ref{def_OmegaM} and Lemma \ref{lem_uniform_multiplier_concentration}, on $\Omega_{\Xi}$,
	\begin{align*}
		\sup_{(x,y)\in\mathcal N_{\delta_0}}
		\|G_n(x,y)-G_{n,c}(x,y)\|=\mathrm{O}(n^{-1/2+c_{\Xi,1}}),
	\end{align*} 
	and
	\begin{align*}
		\sup_{(x,y)\in\mathcal N_{\delta_0}}
		\|DG_n(x,y)-DG_{n,c}(x,y)\|=\mathrm{O}(n^{-1/2+c_{\Xi,1}}).
	\end{align*}
	In particular,
	\begin{align*}
		G_n(x_0,y_0)=G_n(x_0,y_0)-G_{n,c}(x_0,y_0)=\mathrm{O}(n^{-1/2+c_{\Xi,1}}),
	\end{align*}
	and
	\begin{align*}
		DG_n(x_0,y_0)= DG_{n,c}(x_0,y_0)+\mathrm{O}(n^{-1/2+c_{\Xi,1}}).
	\end{align*}
	Hence, for all sufficiently large $n$, $DG_n(x_0,y_0)$ is also invertible and
	\begin{align*}
		\|[DG_n(x_0,y_0)]^{-1}\|\le C
	\end{align*}
	for some constant $C>0$. We now solve $G_n(x,y)=0$ near $(x_0,y_0)$. Write
	\begin{align*}
		d:=
		\begin{pmatrix}
			x-x_0\\
			y-y_0
		\end{pmatrix},
		\qquad
		d_0:=
		-[DG_n(x_0,y_0)]^{-1}G_n(x_0,y_0).
	\end{align*}
	Then $\|d_0\|=\mathrm{O}(n^{-1/2+c_{\Xi,1}})$.
	For $\|d\|\le C_1n^{-1/2+c_{\Xi,1}}$ for some $C_1>0$, Taylor's expansion gives
	\begin{align*}
		G_n(x,y)=G_n(x_0,y_0)+DG_n(x_0,y_0)d+R_n(d),
	\end{align*}
	where, after possibly shrinking $\delta_0$, the second derivatives of $G_n$ are uniformly bounded on $\mathcal N_{\delta_0}$. Hence, we have $ \|R_n(d)\|\le C_2\|d\|^2,$ for some constant $C_2>0$. Equivalently, the equation $G_n(x,y)=0$ can be written as
	\begin{align*}
		d=-[DG_n(x_0,y_0)]^{-1}G_n(x_0,y_0)-
		[DG_n(x_0,y_0)]^{-1}R_n(d)
		=
		d_0-\,[DG_n(x_0,y_0)]^{-1}R_n(d).
	\end{align*}
	Define the map
	\begin{align*}
		\mathcal T(d)
		:=
		d_0-\,[DG_n(x_0,y_0)]^{-1}R_n(d).
	\end{align*}
	For $C_3>C$ large enough and $n$ sufficiently large, $\mathcal T$ maps the
	ball $\mathcal B_n:=\{d:\|d\|\le C_3n^{-1/2+c_{\Xi,1}}$
	into itself, because
	\begin{align*}
		\|\mathcal T(d)\|
		\le
		C n^{-1/2+c_{\Xi,1}}+CC_2\|d\|^2
		\le
		C n^{-1/2+c_{\Xi,1}}+C C_1^2 C_2 n^{-1+2c_{\Xi,1}}
		\le
		C_3n^{-1/2+c_{\Xi,1}}.
	\end{align*}
	Moreover, for $d,d^{\prime}\in\mathcal B_n$,
	\begin{align*}
		\|\mathcal T(d)-\mathcal T(d^{\prime})\|
		\le
		C(\|d\|+\|d^{\prime}\|)\|d-d^{\prime}\|
		\le
		C n^{-1/2+c_{\Xi,1}}\|d-d^{\prime}\|.
	\end{align*}
	Since the coefficient on the right-hand side is $\mathrm{o}(1)$, $\mathcal T$ is a contraction on $\mathcal B_n$ for all sufficiently large $n$. By the contraction mapping theorem, there exists a unique $d_*\in\mathcal B_n$ such that $\mathcal T(d_*)=d_*$. Setting $x_1:=x_0+d_{*,1},
	y_1:=y_0+d_{*,2},$ we obtain $G_n(x_1,y_1)=0;$ that is,
	\begin{align*}
		F_n(x_1,y_1)=0,
		\qquad
		\partial_xF_n(x_1,y_1)=0.
	\end{align*}
	Furthermore,
	\begin{align*}
		|x_1-x_0|+|y_1-y_0|
		\le C\|d_*\|
		=\mathrm{O}(n^{-1/2+c_{\Xi,1}}).
	\end{align*}
	
	The same contraction argument also gives uniqueness in the local ball $\mathcal B_n$. After shrinking $\delta_0$, this is the unique local solution in $\mathcal N_{\delta_0}$. The proof is complete.
\end{proof}

\end{document}
